% sec3a.tex — c06｜原文 p19–22（p19 从 §3 标题后起）｜§3 N=2 理论、§3.1 超对称极限、§3.2 S 对偶、§3.3 开头
% 编号公式 (3.1)–(3.19) 共 19 条；编号环境：习题 3.1、3.2。
% 注意：p19 顶端两句（「这两个条件把变量数减少 2」）已在 sec2b 末尾，本块不重复。
% [原文 p.19]
$\sN=2$ 超共形代数由两个超荷
$\{\mathcal{Q}_{I\alpha},\mathcal{S}^{I\alpha}\equiv\mathcal{Q}^{\dagger I\alpha},\tilde{\mathcal{Q}}^I_{\dot\alpha},\tilde{\mathcal{S}}^{\dot\alpha}_I\equiv\tilde{\mathcal{Q}}^{\dagger\dot\alpha}_I\}$
组成，其中 $I=1,2$。除共形代数外，玻色子代数还包括 R 对称性 $\SU(2)_R\times\U(1)_r$。超荷 $\mathcal{Q}_{I\alpha}$ 在 $\SU(2)$ 下按二重态变换，该 $\SU(2)$ 由指标 $I$ 显式标出。在 $\U(1)_r$ 下，$\mathcal{Q}_{I\alpha}$ 的荷为 $-1$，而 $\tilde{\mathcal{Q}}^I_{\dot\alpha}$ 的荷为 $1$。指标 $\alpha$ 与 $\dot\alpha$ 分别是 $\SU(2)_1$ 与 $\SU(2)_2$ 的指标，而 $\SU(2)_1\times\SU(2)_2=\mathrm{Spin}(4)$ 是共形群的转动子群。所有 $\mathcal{Q}$ 与 $\tilde{\mathcal{Q}}$ 的共形维数都是 $1/2$，所有 $\mathcal{S}$ 与 $\tilde{\mathcal{S}}$ 的共形维数都是 $-1/2$。超共形代数中最重要的（反）对易关系为
\[
\{\mathcal{Q}_{I\alpha},\mathcal{Q}^{\dagger J\beta}\}
=\Delta\delta_\alpha^\beta\delta_I^J+2M_\alpha^\beta\delta_I^J-2R_I^J\delta_\alpha^\beta+\frac12 r\,\delta_\alpha^\beta\delta_I^J ,
\tag{3.1}
\]
\[
\{\tilde{\mathcal{Q}}^I_{\dot\alpha},\tilde{\mathcal{Q}}^{\dagger\dot\beta}_J\}
=\Delta\delta_{\dot\alpha}^{\dot\beta}\delta_J^I+2\tilde M_{\dot\alpha}^{\dot\beta}\delta_J^I-2R_J^I\delta_{\dot\alpha}^{\dot\beta}-\frac12 r\,\delta_{\dot\alpha}^{\dot\beta}\delta_J^I .
\tag{3.2}
\]
不失一般性，我们取 $\mathcal{Q}_{1-}$ 来计算指标。显式地写出下式是有用的：
\[
\{\mathcal{Q}_{1-},(\mathcal{Q}_{1-})^\dagger\}=\Delta-2j_1-2R+\frac12 r .
\tag{3.3}
\]
选定了超荷 $\mathcal{Q}_{1-}$，就确定了一个 $\sN=1$ 子代数。它由
$\{\mathcal{Q}_{1\alpha},\mathcal{S}^{1\alpha}\equiv\mathcal{Q}^{\dagger1\alpha},\tilde{\mathcal{Q}}^1_{\dot\alpha},\tilde{\mathcal{S}}^{\dot\alpha}_1\equiv\tilde{\mathcal{Q}}^{\dagger\dot\alpha}_1\}$
生成。与式 (2.2) 比较，
\[
r_{\sN=1}=\frac23\Big(-2R+\frac12 r\Big) .
\tag{3.4}
\]
此外，所有超荷都与 $R+\frac12 r$ 对易。由于只有超荷与 R 对称性有非平凡的对易关系，$R+\frac12 r$ 就是相对于整个 $\sN=1$ 代数的对易子。我们通过开启关于该荷的逸度来定义指标，
\[
\sI(p,q,t)\equiv\Tr(-1)^F p^{\frac13(\Delta+j_1)+j_2}q^{\frac13(\Delta+j_1)-j_2}x^{(R+\frac12 r)},\qquad
\delta=\Delta-2j_1-2R+\frac12 r ,
\tag{3.5}
\]
\paragraph{超多重态的指标} $\sN=2$ 超多重态由一对手征多重态组成。一个手征多重态与另一个手征多重态的复共轭一起，在 $\SU(2)_R$ 对称性下按二重态 $q_I$ 变换，并且在整体对称性下具有相同的表示。这意味着超多重态中的两个手征多重态按整体对称性的共轭表示变换。我们来看它们对指标的贡献。

% [原文 p.20]
手征多重态 $q_1$ 的 $r_{\sN=1}=\frac23(1+0)$，另一个手征多重态 $\bar q_2$ 也是如此。此外，它们在 $R+\frac12 r=-\frac12$ 这个对易子下的荷相同。因此超多重态的单字母指标为
\[
\begin{aligned}
i_H(a_i;p,q,x)&=\frac{(pq)^{\frac13}x^{\frac12}\chi_R(a_i)-(pq)^{\frac23}x^{-\frac12}\chi_{\bar R}(a_i)}{(1-p)(1-q)}
+\frac{(pq)^{\frac13}x^{\frac12}\chi_{\bar R}(a_i)-(pq)^{\frac23}x^{-\frac12}\chi_R(a_i)}{(1-p)(1-q)}\\
&=\frac{(pq)^{\frac13}x^{\frac12}-(pq)^{\frac23}x^{-\frac12}}{(1-p)(1-q)}\Big(\chi_R(a_i)+\chi_{\bar R}(a_i)\Big) .
\end{aligned}
\tag{3.6}
\]
单个手征多重态 $q_1$ 有时称为半超多重态。超多重态的贡献可以理解为半超多重态取遍表示 $R$ 与 $\bar R$ 时的贡献。

\paragraph{矢量多重态的指标} $\sN=2$ 矢量多重态由一个 $\sN=1$ 手征多重态 $\phi$ 与一个 $\sN=1$ 矢量多重态 $V$ 组成。该手征多重态在 $\SU(2)_R$ 下中性，但在 $\U(1)_r$ 下的荷为 2。有效的 $r_{\sN=1}=\frac23(0+\frac12)=\frac23$，在对易子下的荷为 $R+\frac12 r=1$。当然，$\sN=1$ 矢量多重态不能带有任何整体对称性的荷。因此，$\sN=2$ 矢量多重态的指标为
\[
i_V(a_i;p,q,x)=\Big(\frac{(pq)^{\frac13}x^{-1}-(pq)^{\frac23}x}{(1-p)(1-q)}+\frac{-p-q+2pq}{(1-p)(1-q)}\Big)\chi_{\adj}(a_i) .
\tag{3.7}
\]
\subsection{超对称极限}
可以取超共形指标的极限，使得并非全部短多重态、而只有某些特殊的短多重态对指标有贡献。$\sN=2$ 指标的这类超对称极限已在 [30] 中分类并命名。为便于分析，采用变量 $t$，作替换 $x\to(pq)^{-\frac23}t$。用这些变量，指标为
\[
\sI(p,q,t)=\Tr(-1)^F p^{\frac13(\Delta+j_1)+j_2-\frac23(R+\frac12 r)}q^{\frac13(\Delta+j_1)-j_2-\frac23(R+\frac12 r)}t^{R+\frac12 r} .
\tag{3.8}
\]
利用 $\delta=\Delta-2j_1-2R+\frac12 r=0$，它可以改写为
\[
\begin{aligned}
\sI(p,q,t)&=\Tr(-1)^F p^{\frac12(\Delta+2j_2-2R-\frac12 r)}q^{\frac12(\Delta-2j_2-2R-\frac12 r)}t^{R+\frac12 r},\\
&=\Tr(-1)^F p^{\frac12\{\tilde{\mathcal{Q}}_{1\dot+},(\tilde{\mathcal{Q}}_{1\dot+})^\dagger\}}q^{\frac12\{\tilde{\mathcal{Q}}_{1\dot-},(\tilde{\mathcal{Q}}_{1\dot-})^\dagger\}}t^{R+\frac12 r} .
\end{aligned}
\tag{3.9}
\]
于是半超多重态与矢量多重态的指标变为
\[
i_{\frac12 H}(p,q,t)=\frac{\sqrt t-pq/\sqrt t}{(1-p)(1-q)},\qquad
i_V(p,q,t)=\frac{pq/t-t}{(1-p)(1-q)}+\frac{-p-q+2pq}{(1-p)(1-q)} .
\tag{3.10}
\]
把指标写成式 (3.9) 的形式，其好处是出现的荷显为非负。这允许把 $p$ 或 $q$ 或二者同时取为零。如下文所见，这些极限计入具有增强超对称性的算符。
% [原文 p.21]
\begin{exercise}[3.1]
定义 $\delta_\alpha=\{Q_\alpha,(Q_\alpha)^\dagger\}$，并相应地定义 $\tilde\delta_{\dot\alpha}$，用荷 $\delta_\alpha,\tilde\delta_{\dot\alpha}$ 构造指标。这样选择荷的好处是它们是正定的，因而相应的逸度可以取为零。
\end{exercise}
\paragraph{Macdonald 指标：} 对应 $p\to0$ 的极限。在此极限下，$\{\tilde{\mathcal{Q}}_{1\dot+},(\tilde{\mathcal{Q}}_{1\dot+})^\dagger\}=0$ 的态有贡献。这意味着它们除了被 $\mathcal{Q}_{1-}$$\,$熄灭外，还被 $\tilde{\mathcal{Q}}_{1\dot+}$ 熄灭。在此极限下单字母指标变为
\[
i_{\frac12 H}(q,t)=\frac{\sqrt t}{1-q},\qquad i_V(q,t)=\frac{-t-q}{1-q} .
\tag{3.11}
\]
\paragraph{Hall–Littlewood 指标：} 这是指标的 $p\to0,\ q\to0$ 极限。在此极限下，算符除了被 $\mathcal{Q}_{1-}$ 熄灭外，还被 $\tilde{\mathcal{Q}}_{1\dot+}$ 与 $\tilde{\mathcal{Q}}_{1\dot-}$ 熄灭。条件 $\delta=0$、$\Delta+2j_2-2R-\frac12 r=0$、$\Delta-2j_2-2R-\frac12 r=0$ 可以化简为
\[
j_2=0,\qquad j_1=r,\qquad \Delta=2R+r .
\tag{3.12}
\]
此极限下的单字母指标为
\[
i_{\frac12 H}(t)=\sqrt t,\qquad i_V(t)=-t .
\tag{3.13}
\]
\paragraph{Schur 指标：} Schur 极限是指标的 $q=t$ 极限。
\[
\sI(p,q)=\Tr(-1)^F p^{\frac12(\Delta+2j_2-2R-\frac12 r)}q^{\frac12(\Delta-2j_2-2R-\frac12 r)+(R+\frac12 r)}
\tag{3.14}
\]
指标中出现的荷是 $\Delta+2j_2-2R-\frac12 r$ 与 $\Delta-2j_2+\frac12 r$。由构造，这些荷与 $\mathcal{Q}_{1-}$ 对易，因为我们计算的正是关于 $\mathcal{Q}_{1-}$ 的指标。有趣的是，这些荷也与 $\tilde{\mathcal{Q}}_{1\dot+}$ 对易。因此，指标只从 $\{\tilde{\mathcal{Q}}_{1\dot+},(\tilde{\mathcal{Q}}_{1\dot+})^\dagger\}=0$ 的态获得贡献。于是指标自动与 $p$ 无关，迹公式化简为
\[
\sI(q)=\Tr(-1)^F q^{\Delta-R} .
\tag{3.15}
\]
此极限下的单字母指标为
\[
i_{\frac12 H}(q)=\frac{\sqrt q}{1-q},\qquad i_V(q)=\frac{-2q}{1-q} .
\tag{3.16}
\]
超对称极限使人得以探测 $\sN=2$ 超共形代数的额外超对称扇区。
% [原文 p.22]
\subsection{$\sN=2$ 超 QCD 的 S 对偶}
一大类 $\sN=2$ 规范理论（称为 class S）的非平凡对偶由 Gaiotto 发现 [31]。这类对偶中最简单的一个——已由 Seiberg 与 Witten 的工作 [32] 所知——是带八个半超多重态的 $\SU(2)$ 规范理论。这些半超多重态按 $\SO(8)$ 的矢量表示变换。我们只看子群 $\SO(8)\supset\SO(4)\times\SO(4)=\SU(2)_a\times\SU(2)_b\times\SU(2)_c\times\SU(2)_d$。$\SO(8)$ 按 $\SU(2)^4$ 的表示分解为 $\mathbf{8}_v=(2_a\otimes2_b)\oplus(2_c\otimes2_d)$。该表示的特征标为
\[
\Big(a+\frac1a\Big)\Big(b+\frac1b\Big)+\Big(c+\frac1c\Big)\Big(d+\frac1d\Big) .
\tag{3.17}
\]
有了这些信息，我们立即可以写下指标的积分公式。该理论的指标为
\[
\sI=\kappa\,\Gamma\Big(\frac{pq}{t}\Big)\oint\frac{\dd z}{4\pi\ii z}\frac{\Gamma\big(\frac{pq}{t}z^{\pm2}\big)}{\Gamma(z^{\pm2})}\Gamma\big(\sqrt t\,z^{\pm1}a^{\pm1}b^{\pm1}\big)\Gamma\big(\sqrt t\,z^{\pm1}c^{\pm1}d^{\pm1}\big) .
\tag{3.18}
\]
这里我们用了简写 $\Gamma(z^\pm)=\Gamma(z)\Gamma(z^{-1})$ 等等。

已知该理论具有三重性（triality）。若把半超多重态的矢量表示换成旋量表示或共轭旋量表示，则得到同一理论，但耦合不同。由于指标对耦合不敏感，我们预期指标在这些表示改变下不变。旋量与共轭旋量按 $\SU(2)^4$ 分解为
\[
8_s=(2_a\otimes2_c)\oplus(2_b\otimes2_d),\qquad
8_c=(2_a\otimes2_d)\oplus(2_b\otimes2_c) .
\tag{3.19}
\]
这意味着指标必须在交换 $b\leftrightarrow c$ 与 $b\leftrightarrow d$ 下不变。除这些非平凡的不变性外，指标显然在平凡交换 $a\leftrightarrow b$ 与 $c\leftrightarrow d$ 下不变。

在数学文献中，这些不变性被发现 [33] 的时间与物理学中预期它们 [34] 的时间大致相同。作为物理学家，我们可以通过令 $p=xy$、$q=x/y$ 并按 $x$ 的幂展开来快速检验这些不变性。
\begin{exercise}[3.2]
完成此事。
\end{exercise}
\subsection{class S 对偶与 TQFT}
根据 Gaiotto 的工作，已知该理论是通过把六维秩 1 的 $(2,0)$ 理论紧化在带四个穿刺点的球面上得到的。参数化该曲面的复结构映射为复化规范耦合。详细讨论见 Bruno LeFloch 的讲义。此外，$\SU(2)^4$ 整体对称性中的每个 $\SU(2)$ 因子对应一个穿刺点。

% [Tutor Note: 术语：本节超荷用花体 $\mathcal{Q}$（原文如此），而习题 3.1 与 §2 用正体 $Q$；不要因字形不同而改符号。]
% [Tutor Note: 易错：(3.1) 与 (3.2) 中 $R$ 的指标次序不同（$R_I^J$ 与 $R_J^I$），且 $r$ 项符号相反（$+\frac12r$ 与 $-\frac12r$），这是 $\U(1)_r$ 手征性的体现。]
% [Tutor Note: 易错：本节 $\delta$ 含 $-2R$ 项，与 §2 的 $\delta=\Delta-2j_1+\frac32r$ 不同；换了超荷，$\delta$ 随之改变。]
% [Tutor Note: 动机：(3.4) 把 $\sN=2$ 的 R 荷翻译成 $\sN=1$ 的 $r$ 荷，这样才能对每个 $\sN=1$ 多重态套用 §2 的单字母公式。]
% [Tutor Note: 术语：hypermultiplet 译「超多重态」，half-hypermultiplet 译「半超多重态」；二者都不是 $\sN=1$ 的独立多重态。]
% [Tutor Note: 易错：(3.6) 第一行两个分式中 $\chi_R$ 与 $\chi_{\bar R}$ 的次序相反，第二个分式来自 $\bar q_2$；不要抄成同一个分式。]
% [Tutor Note: 动机：(3.9) 把指数写成超荷反对易子的形式，使荷显为非负，这是取 $p,q\to0$ 极限（Macdonald/Hall–Littlewood）的前提。]
% [Tutor Note: 易错：(3.9) 中 $p,q$ 的指数分别是两个反对易子的一半，不是 (3.5) 的 $\frac13(\Delta+j_1)\pm j_2$；两者只在 $\delta=0$ 时等价。]
% [Tutor Note: 物理：三个极限分别对应 $p\to0$（Macdonald）、$p,q\to0$（Hall–Littlewood）、$q=t$（Schur），各自只保留增强超对称的扇区。]
% [Tutor Note: 易错：Schur 极限是令 $q=t$，不是令 $t=0$；结果 (3.15) 只依赖一个变量 $q$。]
% [Tutor Note: 物理：(3.12) 的三个条件 $j_2=0$、$j_1=r$、$\Delta=2R+r$ 是 Hall–Littlewood 极限下存活态的 BPS 条件。]
% [Tutor Note: 术语：triality 译「三重性」；$\mathbf{8}_v,\mathbf{8}_s,\mathbf{8}_c$ 的分解正是 $\SO(8)$ 三重性的群论体现。]
% [Tutor Note: 易错：(3.18) 中被积函数含两个 $\Gamma$ 因子（分别含 $a,b$ 与 $c,d$），不是四个；$z^{\pm2}$ 是 $z^2,z^{-2}$ 的简写。]
% [Tutor Note: 文献：[31] Gaiotto 的 class S 构造；[32] Seiberg–Witten；[33][34] 数学与物理几乎同时发现这些不变性。]
% [Tutor Note: 物理：指标对耦合不敏感，因此 S 对偶两边指标必须相等；这是检验对偶的强约束，但不是对偶的证明。]
% [Tutor Note: 语病：原文 p.22 「the a puncture」冠词衍余，应为「a puncture」；照译并登记 ERRATA。]
